\section{Quantum shuf\/f\/le characters}\label{sec:shuffle_character}

The results of Sections~\ref{sec:feigin},~\ref{subsec:shuffle_gen_feigin}, and~\ref{subsec:hall_gen_feigin} have
established the commutativity of the front two faces in the tetrahedron below:
\begin{gather}\label{dia:feigin_tetrahedron}
\begin{split}&
\includegraphics{Rupel-Fig1}
\end{split}
\end{gather}
%\begin{equation}%
%\label{dia:feigin_tetrahedron}
% \begin{tikzpicture}[description/.style={fill=white,inner sep=2pt}]
%     \matrix (m) [matrix of math nodes, row sep=3em,column sep=2.5em, text height=1.5ex, text depth=0.25ex]
%      { & \mathcal{H}^*(Q,\mathbf{d}) & & \\ \mathcal{A}_v[N] & & & P_\mathbf{i} \\ & & g^* & \\};
%     \path[->,font=\scriptsize]
%      (m-1-2) edge[dashed] node[near start, above=4pt] {$\Omega$} (m-3-3)
%      (m-2-1) edge (m-1-2) edge (m-3-3) edge[line width=6pt,draw=white] (m-2-4) edge node[near start,above] {$\Psi_\mathbf{i}$} (m-2-4)
%      (m-1-2) edge node[above] {$\widetilde\Psi_\mathbf{i}$} (m-2-4)
%      (m-3-3) edge node[below=3pt] {$\overline\Psi_\mathbf{i}$} (m-2-4);
%\end{tikzpicture}
%\end{equation}
Our goal of this section is to complete this diagram, that is def\/ine the \emph{quantum shuffle character} homomorphism
$\Omega: \mathcal{H}^*(Q,\mathbf{d})\to g^*$ making the remaining two faces of the tetrahedron commute.

For a~word $\mathbf{j}=(j_1,\ldots,j_r)$ write $\mathcal{F}_\mathbf{j}(V)$ for the set of f\/lags in~$V$ of type~$\mathbf{j}$:
\begin{gather*}
\mathcal{F}_\mathbf{j}(V)=\{0=V_r\subset V_{r-1}\subset\dots\subset V_1\subset V_0=V: V_{k-1}/V_k\cong S_{j_k}\;\text{for $1\le k\le r$}\}. %\;
\end{gather*}
Then for a~basis vector $[V]^*\in\mathcal{H}^*(Q,\mathbf{d})$ we def\/ine $\Omega([V]^*)\in g^*$ as a~certain generating
function for counting f\/lags in~$V$:
\begin{gather}
\label{eq:Omega}
\Omega([V]^*)=\sum\limits_{\mathbf{j}\in\mathbf{W}} v^{-\sum\limits_{k<\ell}\langle\alpha_{j_\ell},\alpha_{j_k}\rangle}
|\mathcal{F}_\mathbf{j}(V)|\cdot\mathbf{j}.
\end{gather}
This map $\Omega:\mathcal{H}^*(Q,\mathbf{d})\to g^*$ turns out to preserve a~fair amount of structure.

\begin{Theorem}
\label{th:shuffle_character}
The map $\Omega: \mathcal{H}^*(Q,\mathbf{d})\to g^*$ given by~\eqref{eq:Omega} is a~homomorphism of twisted bialgebras.
\end{Theorem}

\begin{Remark}
In general,~$\Omega$ will not be injective.
For example, every member of the $\mathbb{P}^1$-family of indecomposable representations with dimension vector $(1,1)$
for the 2-Kronecker quiver will have the same image under~$\Omega$.
\end{Remark}

\begin{proof}
We begin by showing that~$\Omega$ is an algebra homomorphism.
Recall the rescaling
\begin{gather*}
[V]^*=v^{-\frac{1}{2}\langle |V|,|V|\rangle+\frac{1}{2}\sum\limits_{i=1}^n d_iv_i}\delta_{[V]}.
\end{gather*}
The product $[U]^*[W]^*$ can then be computed in this same basis as
\begin{gather*}
[U]^*[W]^*=v^{-\frac{1}{2}\langle |U|+|W|,|U|+|W|\rangle+\frac{1}{2}\langle |U|,|W|\rangle+\frac{1}{2}\langle
|W|,|U|\rangle+\frac{1}{2}\sum\limits_{i=1}^n d_i(u_i+w_i)}\delta_{[U]}\delta_{[W]}
\\
\phantom{[U]^*[W]^*}
 =\sum\limits_{[V]}v^{-\frac{1}{2}\langle |V|,|V|\rangle+\frac{1}{2}\langle |U|,|W|\rangle+\frac{1}{2}\langle
|W|,|U|\rangle+\frac{1}{2}\sum\limits_{i=1}^n d_iv_i}v^{\langle|U|,|W|\rangle}
\\
\phantom{[U]^*[W]^*=}{}
\times
\frac{|\operatorname{Aut}(U)||\operatorname{Aut}(W)|}{|\operatorname{Aut}(V)|}\cdot\big|\mathcal{F}^V_{UW}\big|\cdot \delta_{[V]}
\\
%\label{eq:hall_basis_prod}
\phantom{[U]^*[W]^*}
 =\sum\limits_{[V]}v^{\frac{3}{2}\langle |U|,|W|\rangle+\frac{1}{2}\langle
|W|,|U|\rangle}\frac{|\operatorname{Aut}(U)||\operatorname{Aut}(W)|}{|\operatorname{Aut}(V)|}\cdot\big|\mathcal{F}^V_{UW}\big|\cdot [V]^*.
\end{gather*}
Thus the equality $\Omega([U]^*[W]^*)=\Omega([U]^*)\Omega([W]^*)$ is equivalent to the following interesting collection
of identities involving f\/lags.
\begin{Lemma}
\label{lem:flag_identity}
For every word $\mathbf{h}\in\mathbf{W}$ we have
\begin{gather}
\nonumber
\sum\limits_{[V]} v^{\frac{3}{2}\langle |U|,|W|\rangle+\frac{1}{2}\langle|W|,|U|\rangle-\sum\limits_{k<\ell}\langle\alpha_{h_\ell},\alpha_{h_k}\rangle}
\frac{|\operatorname{Aut}(U)||\operatorname{Aut}(W)|}{|\operatorname{Aut}(V)|}\big|\mathcal{F}^V_{UW}\big||\mathcal{F}_\mathbf{h}(V)|
\\
\qquad
 =\sum\limits_{\substack{\mathbf{j}_1,\mathbf{j}_2\in\mathbf{W}\\
\sigma\in\Sigma_{\operatorname{ht}(U),\operatorname{ht}(W)}: \mathbf{j}_\sigma=\mathbf{h}}}v^{-\sum\limits_{1\le
k<\ell\le t}\langle\alpha_{j_\ell},\alpha_{j_k}\rangle-\sum\limits_{t+1\le k<\ell\le
r+s}\langle\alpha_{j_\ell},\alpha_{j_k}\rangle+\zeta(\sigma)}|\mathcal{F}_{\mathbf{j}_1}(U)||\mathcal{F}_{\mathbf{j}_2}(W)|,
\label{eq:flag_identity}
\end{gather}
where we abbreviated $t=\operatorname{ht}(U)$.
\end{Lemma}

\begin{proof}
This essentially follows from the compatibility of multiplication and comultiplication in the Hall--Ringel algebra
$\mathcal{H}(Q,\mathbf{d})$.
Using the def\/initions it is easy to compute the following composition of comultiplication and iterated multiplication:
\begin{gather*}
\Delta([S_{h_1}]\cdots[S_{h_{r+s}}])=\sum\limits_{[V]}v^{\sum\limits_{k<\ell}\langle
\alpha_{h_k},\alpha_{h_\ell}\rangle}|\mathcal{F}_\mathbf{h}(V)|\Delta([V])
\\
%\label{eq:comulit_iter_mult}
\qquad
 =\sum\limits_{[U],[V],[W]}v^{\langle|U|,|W|\rangle+\sum\limits_{k<\ell}\langle\alpha_{h_k},\alpha_{h_\ell}\rangle}
\frac{|\operatorname{Aut}(U)||\operatorname{Aut}(W)|}{|\operatorname{Aut}(V)|}\big|\mathcal{F}^V_{UW}\big||\mathcal{F}_\mathbf{h}(V)|\cdot[U]\otimes[V].
\end{gather*}
Theorem~\ref{th:hall_comult} asserts that this is the same as the following iterated multiplication of
comultiplications:
\begin{gather*}
\Delta([S_{h_1}])\cdots\Delta([S_{h_{r+s}}])=([S_{h_1}]\otimes1+1\otimes[S_{h_1}])\cdots([S_{h_{r+s}}]\otimes1+1\otimes[S_{h_{r+s}}])
\\
%\label{eq:iter_mult_comult}
\qquad
 =\sum\limits_{\substack{[U],[W]:\\ \operatorname{ht}(U)+\operatorname{ht}(W)=r+s\\
\sigma\in\Sigma_{\operatorname{ht}(U),\operatorname{ht}(W)}\\
\mathbf{j}_1,\mathbf{j}_2\in\mathbf{W}: \mathbf{j}_\sigma=\mathbf{h}}}v^E
|\mathcal{F}_{\mathbf{j}_1}(U)||\mathcal{F}_{\mathbf{j}_2}(W)|\cdot [U]\otimes[W],
\end{gather*}
with exponent $E=\sum\limits_{1\le k<\ell\le t}\langle\alpha_{j_k},\alpha_{j_\ell}\rangle
+\sum\limits_{t+1\le k<\ell\le r+s}\langle\alpha_{j_k},\alpha_{j_\ell}\rangle
+\sum\limits_{\substack{1\le k\le t\\t+1\le\ell\le r+s\\\sigma_k^{-1}>\sigma_\ell^{-1}}}(\alpha_{j_k},\alpha_{j_\ell})$,
where we abbreviate $t=\operatorname{ht}(U)$.

A few remarks seem necessary to clarify the second equality above.
To obtain a~f\/ixed $[U]$ and $[W]$ we must take from the product
\begin{gather*}
([S_{h_1}]\otimes1+1\otimes[S_{h_1}])\cdots([S_{h_{r+s}}]\otimes1+1\otimes[S_{h_{r+s}}])
\end{gather*}
exactly $\operatorname{ht}(U)$ factors, recorded by $\mathbf{j}_1$, of the form $[S_{h_k}]\otimes1$ and
$\operatorname{ht}(W)$ factors, recorded by $\mathbf{j}_2$, of the form $1\otimes[S_{h_k}]$.
With this notation we are looking at the products $[S_{j_1}]\cdots[S_{j_t}]\otimes[S_{j_{t+1}}]\cdots[S_{j_{r+s}}]$, where
again we have used the abbreviation $\operatorname{ht}(U)=t$.
The f\/irst and second sums in~$E$ as well as the f\/lag coef\/f\/icients should now be transparent.
In order to record the contribution of the twisted multiplication on
$\mathcal{H}(Q,\mathbf{d})\otimes\mathcal{H}(Q,\mathbf{d})$ to the coef\/f\/icient of $[U]\otimes[W]$ we need a~shuf\/f\/le
$\sigma\in\Sigma_{\operatorname{ht}(U),\operatorname{ht}(W)}$ and $j_k=h_{\sigma_k^{-1}}$.
Indeed, the condition $\sigma_k^{-1}>\sigma_\ell^{-1}$ of the last sum in~$E$ records a~product of the form
$(1\otimes[S_{j_k}])([S_{j_\ell}]\otimes1)$ and thus produces the stated contribution.

As previously stated, Theorem~\ref{th:hall_comult} allows us, for a~f\/ixed $[U]\otimes[W]$, to extract the equality below
from which the claim will follow:
\begin{gather}
\sum\limits_{[V]}v^{\langle|U|,|W|\rangle+\sum\limits_{k<\ell}\langle\alpha_{h_k},\alpha_{h_\ell}\rangle}
\frac{|\operatorname{Aut}(U)||\operatorname{Aut}(W)|}{|\operatorname{Aut}(V)|}|\mathcal{F}^V_{UW}||\mathcal{F}_\mathbf{h}(V)|
\nonumber
\\
\qquad
=\sum\limits_{\substack{\sigma\in\Sigma_{\operatorname{ht}(U),\operatorname{ht}(W)}\\
\mathbf{j}_1,\mathbf{j}_2\in\mathbf{W}: \mathbf{j}_\sigma=\mathbf{h}}}
v^E|\mathcal{F}_{\mathbf{j}_1}(U)||\mathcal{F}_{\mathbf{j}_2}(W)|.
\label{eq:green_equality}
\end{gather}
To complete the proof we perform a~few simple manipulations with the sums appearing in the exponents
of~\eqref{eq:flag_identity} and~\eqref{eq:green_equality}.
To start we match the left hand sides of~\eqref{eq:flag_identity} and~\eqref{eq:green_equality} by multiplying both
sides of~\eqref{eq:flag_identity} by $v^{\sum\limits_{k<\ell}\langle
\alpha_{h_\ell},\alpha_{h_k}\rangle-\frac{1}{2}\langle |U|,|W|\rangle-\frac{1}{2}\langle |W|,|U|\rangle}$ and both sides
of~\eqref{eq:green_equality} by $v^{-\sum\limits_{k<\ell}\langle \alpha_{h_k},\alpha_{h_\ell}\rangle}$.
Thus it remains to establish the following identity of exponents.
\begin{Lemma}
\label{lem:exp_equality}
For any $\sigma\in\Sigma_{\operatorname{ht}(U),\operatorname{ht}(W)}$ and words $\mathbf{j}_1$, $\mathbf{j}_2$ satisfying
$\mathcal{F}_{\mathbf{j}_1}(U)\ne\varnothing$, $\mathcal{F}_{\mathbf{j}_2}(W)\ne\varnothing$, and
$\mathbf{j}_\sigma=\mathbf{h}$ we have:
\begin{gather}
\nonumber
\sum\limits_{k<\ell}\langle \alpha_{h_\ell},\alpha_{h_k}\rangle
-\!\sum\limits_{1\le k<\ell\le t}\!\!\langle\alpha_{j_\ell},\alpha_{j_k}\rangle
-\!\sum\limits_{t+1\le k<\ell\le r+s}\!\!\langle\alpha_{j_\ell},\alpha_{j_k}\rangle
+\zeta(\sigma)-\frac{1}{2}\langle |U|,|W|\rangle-\frac{1}{2}\langle|W|,|U|\rangle
\nonumber
\\
\qquad
=-\sum\limits_{k<\ell}\langle \alpha_{h_k},\alpha_{h_\ell}\rangle+\sum\limits_{1\le k<\ell\le t}\!\!\langle\alpha_{j_k},\alpha_{j_\ell}\rangle
+\sum\limits_{t+1\le k<\ell\le r+s}\!\!\!\!\langle\alpha_{j_k},\alpha_{j_\ell}\rangle
+\!\!\!\sum\limits_{\substack{1\le k\le t\\t+1\le\ell\le r+s\\ \sigma_k^{-1}>\sigma_\ell^{-1}}}\!\!\!\!(\alpha_{j_k},\alpha_{j_\ell}). %\!\!
\label{eq:exp_equality}
\end{gather}
\end{Lemma}

\begin{proof}
We f\/irst eliminate the dependence on $\mathbf{h}$ using the following identities which are immediate consequences of the
equality $j_k=h_{\sigma_k^{-1}}$:
\begin{gather*}
\sum\limits_{k<\ell}\langle \alpha_{h_\ell},\alpha_{h_k}\rangle
-\sum\limits_{1\le k<\ell\le t}\langle\alpha_{j_\ell},\alpha_{j_k}\rangle
-\sum\limits_{t+1\le k<\ell\le r+s}\langle\alpha_{j_\ell},\alpha_{j_k}\rangle
\\
\qquad
=\sum\limits_{\substack{1\le k\le t\\t+1\le\ell\le r+s\\ \sigma_k^{-1}<\sigma_\ell^{-1}}}
\langle\alpha_{j_\ell},\alpha_{j_k}\rangle
+\sum\limits_{\substack{1\le k\le t\\t+1\le\ell\le r+s\\ \sigma_\ell^{-1}<\sigma_k^{-1}}}
\langle\alpha_{j_k},\alpha_{j_\ell}\rangle,
\\
-\sum\limits_{k<\ell}\langle \alpha_{h_k},\alpha_{h_\ell}\rangle
+\sum\limits_{1\le k<\ell\le t}\langle\alpha_{j_k},\alpha_{j_\ell}\rangle
+\sum\limits_{t+1\le k<\ell\le r+s}\langle\alpha_{j_k},\alpha_{j_\ell}\rangle
\\
\qquad
=-\sum\limits_{\substack{1\le k\le t\\t+1\le\ell\le r+s\\ \sigma_k^{-1}<\sigma_\ell^{-1}}}
\langle\alpha_{j_k},\alpha_{j_\ell}\rangle
-\sum\limits_{\substack{1\le k\le t\\t+1\le\ell\le r+s\\ \sigma_\ell^{-1}<\sigma_k^{-1}}}
\langle\alpha_{j_\ell},\alpha_{j_k}\rangle,
\end{gather*}
this reduces~\eqref{eq:exp_equality} to showing
\begin{gather}
\nonumber
\sum\limits_{\substack{1\le k\le t\\t+1\le\ell\le r+s\\ \sigma_k^{-1}<\sigma_\ell^{-1}}}\langle\alpha_{j_\ell},\alpha_{j_k}\rangle
+\sum\limits_{\substack{1\le k\le t\\t+1\le\ell\le r+s\\ \sigma_\ell^{-1}<\sigma_k^{-1}}}\langle\alpha_{j_k},\alpha_{j_\ell}\rangle
+\zeta(\sigma)-\frac{1}{2}\langle|U|,|W|\rangle-\frac{1}{2}\langle |W|,|U|\rangle
\\
\qquad
=-\sum\limits_{\substack{1\le k\le t\\t+1\le\ell\le r+s\\ \sigma_k^{-1}<\sigma_\ell^{-1}}}\langle\alpha_{j_k},\alpha_{j_\ell}\rangle
-\sum\limits_{\substack{1\le k\le t\\t+1\le\ell\le r+s\\ \sigma_\ell^{-1}<\sigma_k^{-1}}}\langle\alpha_{j_\ell},\alpha_{j_k}\rangle
+\sum\limits_{\substack{1\le k\le t\\t+1\le\ell\le r+s\\ \sigma_k^{-1}>\sigma_\ell^{-1}}}(\alpha_{j_k},\alpha_{j_\ell}).
\label{eq:exp_equality2}
\end{gather}
Using the dichotomy $\sigma_k^{-1}<\sigma_\ell^{-1}$ or $\sigma_\ell^{-1}<\sigma_k^{-1}$ for $1\le k\le t$ and
$t+1\le\ell\le r+s$ the last sum on the right hand side of~\eqref{eq:exp_equality2} becomes:
\begin{gather}
\label{eq:zeta}
\sum\limits_{\substack{1\le k\le t\\t+1\le\ell\le r+s\\ \sigma_k^{-1}>\sigma_\ell^{-1}}}(\alpha_{j_k},\alpha_{j_\ell})
=\zeta(\sigma)+\frac{1}{2}\sum\limits_{\substack{1\le k\le t\\t+1\le\ell\le r+s}}(\alpha_{j_k},\alpha_{j_\ell}).
\end{gather}
Rewriting
$(\alpha_{j_k},\alpha_{j_\ell})=\langle\alpha_{j_k},\alpha_{j_\ell}\rangle+\langle\alpha_{j_\ell},\alpha_{j_k}\rangle$
and applying~\eqref{eq:zeta} in the right hand side of~\eqref{eq:exp_equality2} gives:
\begin{gather*}
-\sum\limits_{\substack{1\le k\le t\\t+1\le\ell\le r+s\\ \sigma_k^{-1}<\sigma_\ell^{-1}}}\langle\alpha_{j_k},\alpha_{j_\ell}\rangle
-\sum\limits_{\substack{1\le k\le t\\t+1\le\ell\le r+s\\ \sigma_\ell^{-1}<\sigma_k^{-1}}}\langle\alpha_{j_\ell},\alpha_{j_k}\rangle
+\zeta(\sigma)+\frac{1}{2}\sum\limits_{\substack{1\le k\le t\\t+1\le\ell\le r+s}}
(\langle\alpha_{j_k},\alpha_{j_\ell}\rangle+\langle\alpha_{j_\ell},\alpha_{j_k}\rangle)
\\
\qquad
=-\frac{1}{2}\sum\limits_{\substack{1\le k\le t\\t+1\le\ell\le r+s}}
(\langle\alpha_{j_k},\alpha_{j_\ell}\rangle+\langle\alpha_{j_\ell},\alpha_{j_k}\rangle)+\zeta(\sigma)
\\
\qquad
\phantom{=}
+\sum\limits_{\substack{1\le k\le t\\t+1\le\ell\le r+s\\ \sigma_k^{-1}<\sigma_\ell^{-1}}}\langle\alpha_{j_\ell},\alpha_{j_k}\rangle
+\sum\limits_{\substack{1\le k\le t\\t+1\le\ell\le r+s\\ \sigma_\ell^{-1}<\sigma_k^{-1}}}\langle\alpha_{j_k},\alpha_{j_\ell}\rangle
\\
\qquad
=-\frac{1}{2}\langle |U|,|W|\rangle-\frac{1}{2}\langle |W|,|U|\rangle+\zeta(\sigma)
+\sum\limits_{\substack{1\le k\le t\\t+1\le\ell\le r+s\\ \sigma_k^{-1}<\sigma_\ell^{-1}}}\langle\alpha_{j_\ell},\alpha_{j_k}\rangle
+\sum\limits_{\substack{1\le k\le t\\t+1\le\ell\le r+s\\ \sigma_\ell^{-1}<\sigma_k^{-1}}}\langle\alpha_{j_k},\alpha_{j_\ell}\rangle
\end{gather*}
as desired, which completes the proof of Lemma~\ref{lem:exp_equality}.
\end{proof}

This completes the proof of Lemma~\ref{lem:flag_identity}.
\end{proof}

With Lemma~\ref{lem:flag_identity} proven, we may conclude that~$\Omega$ def\/ines a~homomorphism of algebras.
It remains to show that $\Delta\Omega([V]^*)=(\Omega\otimes\Omega)(\Delta([V]^*))$.
The coproduct $\Delta([V]^*)$ can be expanded in the same rescaled basis as
\begin{gather*}
\Delta([V]^*)=v^{-\frac{1}{2}\langle |V|,|V|\rangle+\frac{1}{2}\sum\limits_{i=1}^n d_iv_i}\Delta(\delta_{[V]})
\\
\phantom{\Delta([V]^*)}
 =\sum\limits_{[U],[W]}v^{-\frac{1}{2}\langle |U|+|W|,|U|+|W|\rangle+\frac{1}{2}\sum\limits_{i=1}^n
d_i(u_i+w_i)}v^{\langle|U|,|W|\rangle}\big|\mathcal{F}^V_{UW}\big|\cdot \delta_{[U]}\otimes\delta_{[W]}
\\
%\label{eq:hall_basis_coprod}
\phantom{\Delta([V]^*)}
 =\sum\limits_{[V]}v^{\frac{1}{2}\langle |U|,|W|\rangle-\frac{1}{2}\langle |W|,|U|\rangle} \big|\mathcal{F}^V_{UW}\big|\cdot
[U]^*\otimes[W]^*.
\end{gather*}
Then the identity $\Delta\Omega([V]^*)=(\Omega\otimes\Omega)(\Delta([V]^*))$ is equivalent to the following collection
of identities involving f\/lags.
\begin{Lemma}
For all words $\mathbf{j}_1,\mathbf{j}_2\in\mathbf{W}$ we have
\begin{gather*}
 v^{-\sum\limits_{k<\ell}\langle
\alpha_{h_\ell},\alpha_{h_k}\rangle}v^{\frac{1}{2}(|\mathbf{j}_1|,|\mathbf{j}_2|)}|\mathcal{F}_\mathbf{h}(V)|
 =\sum\limits_{[U],[W]}v^{-\sum\limits_{1\le k<\ell\le r}\langle\alpha_{j_\ell},\alpha_{j_k}\rangle}
v^{-\sum\limits_{r+1\le k<\ell\le r+s}\langle\alpha_{j_\ell},\alpha_{j_k}\rangle}
\\
\qquad
\phantom{=}{}
\times
v^{\frac{1}{2}\langle|U|,|W|\rangle-\frac{1}{2}\langle|W|,|U|\rangle}
\big|\mathcal{F}_{U,W}^V\big|\cdot|\mathcal{F}_{\mathbf{j}_1}(U)|\cdot|\mathcal{F}_{\mathbf{j}_2}(W)|,
\end{gather*}
where we used the word $\mathbf{h}=(\mathbf{j}_1,\mathbf{j}_2)$.
\end{Lemma}

\begin{proof}
This essentially follows from the associativity of multiplication in the Hall--Ringel algebra
$\mathcal{H}(Q,\mathbf{d})$.
Indeed, expanding the products
\begin{gather*}
[S_{h_1}]\cdots[S_{h_{r+s}}]=([S_{j_1}]\cdots[S_{j_r}])\cdot([S_{j_{r+1}}]\cdots[S_{j_{r+s}}])
\end{gather*}
gives
$|\mathcal{F}_\mathbf{h}(V)|=|\mathcal{F}_{U,W}^V|\cdot|\mathcal{F}_{\mathbf{j}_1}(U)|\cdot|\mathcal{F}_{\mathbf{j}_2}(W)|$.
To see the equality of the exponents on either side we note that
\begin{gather*}
\frac{1}{2}(\mathbf{j}_1,\mathbf{j}_2)
=\frac{1}{2}\langle|\mathbf{j}_1|,|\mathbf{j}_2|\rangle+\frac{1}{2}\langle|\mathbf{j}_2|,|\mathbf{j}_1|\rangle
=\frac{1}{2}\langle|U|,|W|\rangle+\frac{1}{2}\langle|W|,|U|\rangle.
\end{gather*}
Combining this with the equality
\begin{gather*}
-\sum\limits_{k<\ell}\langle \alpha_{h_\ell},\alpha_{h_k}\rangle=-\sum\limits_{1\le k<\ell\le r}\langle
\alpha_{j_\ell},\alpha_{j_k}\rangle-\sum\limits_{r+1\le k<\ell\le r+s}\langle
\alpha_{j_\ell},\alpha_{j_k}\rangle-\langle|\mathbf{j}_2|,|\mathbf{j}_1|\rangle
\end{gather*}
completes the proof.
\end{proof}

To f\/inish the proof of Theorem~\ref{th:shuffle_character} we note that~$\Omega$ respects the grading and thus, since the
multiplication on both $\mathcal{H}^*(Q,\mathbf{d})\otimes\mathcal{H}^*(Q,\mathbf{d})$ and $g^*\otimes g^*$ are twisted
by the grading, we may conclude that~$\Omega$ is a~homomorphism of twisted bialgebras.
\end{proof}
